\section{Use bounds with a stated numeric domain}\label{maze-p06}

Each direction changes one coordinate by minus one, zero, or one. Since the original coordinate lies between zero and seven, a candidate coordinate lies between minus one and eight. The reference checks it with an unsigned comparison against eight. Values zero through seven pass; eight and the all-one representation of minus one fail.

The candidate coordinate is formed by a word arithmetic instruction, whose result is sign extended in the MARS MIPS64 environment. For the minus-one case, the full register therefore represents a very large unsigned value, so the unsigned less-than test rejects it. This is a domain-specific bounds idiom, not a claim that signed and unsigned comparisons generally mean the same thing.

The alternative implementation checks the negative case explicitly and then performs a signed upper-bound comparison. It passes the same contract and illustrates a more literal way to express the interval. Neither implementation may use a flattened-index test as a substitute for the coordinate checks.
